\documentclass[tikz,border=0pt]{standalone}
\usepackage{fontspec}
\usetikzlibrary{arrows.meta}

\setmainfont{Arial}
\setsansfont{Arial}

\definecolor{pagebg}{HTML}{F8FAFC}
\definecolor{border}{HTML}{CBD5E1}
\definecolor{textmain}{HTML}{0F172A}
\definecolor{textmuted}{HTML}{475569}
\definecolor{blue}{HTML}{2563EB}
\definecolor{bluefill}{HTML}{DBEAFE}
\definecolor{cyanfill}{HTML}{E0F2FE}
\definecolor{cyanline}{HTML}{0284C7}
\definecolor{red}{HTML}{DC2626}
\definecolor{redfill}{HTML}{FEE2E2}
\definecolor{green}{HTML}{15803D}
\definecolor{greenfill}{HTML}{DCFCE7}

\begin{document}
\begin{tikzpicture}[x=1mm,y=1mm,line cap=round,line join=round,font=\sffamily]
  \fill[pagebg] (0,0) rectangle (360,230);
  \draw[fill=white,draw=border,line width=0.55mm,rounded corners=4mm]
    (6,6) rectangle (354,224);

  \node[font=\sffamily\bfseries\fontsize{22}{27}\selectfont,text=textmain]
    at (180,208) {Давление на глубине};
  \node[font=\sffamily\fontsize{14}{17}\selectfont,text=textmuted]
    at (180,192) {Полное давление складывается из давления на поверхности и веса столба жидкости};

  % Сосуд и покоящаяся жидкость
  \draw[fill=pagebg,draw=border,line width=0.45mm,rounded corners=2.5mm]
    (14,45) rectangle (238,180);
  \fill[cyanfill] (32,58) rectangle (206,145);
  \draw[draw=textmuted,line width=0.8mm]
    (32,165) -- (32,58) -- (206,58) -- (206,165);
  \draw[draw=cyanline,line width=0.8mm] (32,145) -- (206,145);

  \node[font=\sffamily\bfseries\fontsize{14}{17}\selectfont,text=red]
    at (119,171) {давление $p_0$ на поверхности};
  \foreach \x in {59,99,139,179} {
    \draw[-{Stealth[length=3.4mm,width=2.5mm]},draw=red,line width=1.0mm]
      (\x,164) -- (\x,151);
  }

  % Мысленно выделенный столб жидкости над точкой A
  \draw[fill=bluefill,draw=blue,line width=0.55mm,dashed,rounded corners=1.5mm]
    (103,82) rectangle (145,145);
  \node[align=center,font=\sffamily\bfseries\fontsize{13.5}{16.5}\selectfont,text=blue]
    at (124,116) {столб\\жидкости};
  \node[align=center,font=\sffamily\fontsize{13}{16}\selectfont,text=textmuted]
    at (73,102) {его вес создаёт\\добавку $\rho gh$};

  \fill[blue] (124,82) circle (2.2mm);
  \draw[draw=blue,line width=0.55mm] (128,85) -- (146,93);
  \node[anchor=west,fill=white,rounded corners=1.3mm,inner xsep=2.2mm,inner ysep=1.3mm,
    font=\sffamily\bfseries\fontsize{13.5}{16.5}\selectfont,text=blue]
    at (147,94) {точка $A$};

  % Глубина h отсчитывается вертикально от поверхности
  \draw[draw=textmuted,line width=0.45mm,dashed] (145,145) -- (225,145);
  \draw[draw=textmuted,line width=0.45mm,dashed] (145,82) -- (225,82);
  \draw[{Stealth[length=3.2mm,width=2.3mm]}-{Stealth[length=3.2mm,width=2.3mm]},
    draw=green,line width=0.9mm] (219,145) -- (219,82);
  \node[fill=pagebg,inner sep=1.3mm,font=\sffamily\bfseries\fontsize{16}{19}\selectfont,text=green]
    at (219,113.5) {$h$};

  % Смысл двух слагаемых
  \draw[fill=white,draw=border,line width=0.45mm,rounded corners=2.5mm]
    (248,45) rectangle (346,180);

  \draw[fill=bluefill,draw=blue!35,line width=0.4mm,rounded corners=1.8mm]
    (256,139) rectangle (338,171);
  \node[font=\sffamily\bfseries\fontsize{13.5}{16.5}\selectfont,text=blue]
    at (297,162) {давление в точке $A$};
  \node[font=\sffamily\bfseries\fontsize{19}{23}\selectfont,text=textmain]
    at (297,149) {$p_A=p_0+\rho gh$};

  \draw[fill=redfill,draw=red!30,line width=0.4mm,rounded corners=1.8mm]
    (256,99) rectangle (338,129);
  \node[font=\sffamily\bfseries\fontsize{15}{18}\selectfont,text=red]
    at (272,114) {$p_0$};
  \node[anchor=west,align=left,font=\sffamily\fontsize{13}{16}\selectfont,text=textmain]
    at (284,114) {вклад давления\\на поверхности};

  \draw[fill=greenfill,draw=green!35,line width=0.4mm,rounded corners=1.8mm]
    (256,59) rectangle (338,89);
  \node[font=\sffamily\bfseries\fontsize{15}{18}\selectfont,text=green]
    at (275,74) {$\rho gh$};
  \node[anchor=west,align=left,font=\sffamily\fontsize{13}{16}\selectfont,text=textmain]
    at (290,74) {вклад веса\\жидкости};

  \draw[fill=cyanfill,draw=cyanline!45,line width=0.45mm,rounded corners=2mm]
    (14,14) rectangle (346,36);
  \node[font=\sffamily\bfseries\fontsize{13.5}{16.5}\selectfont,text=cyanline]
    at (180,25) {Глубину $h$ отсчитывают вертикально вниз от поверхности жидкости};
\end{tikzpicture}
\end{document}
